% MT021 / explanation — canonical editable content.
\(f(x)=w^Tx+b\)只说明预测时怎样算分数：各特征乘权重后相加，再加偏置。它没有规定训练时怎样选\(w,b\)。同一批点可以画出很多分隔线，不同模型用不同标准选择其中一条。


\textbf{逻辑回归（LR）看真实标签有多可信。} 分数经sigmoid变成类别概率。如果真实标签是1，模型给它0.9比给它0.6更好。训练把所有样本的真实标签概率相乘并尽量增大；常用等价的负log损失来最小化。正则化再控制权重大小。


\textbf{SVM看间隔及违例。} 它希望分隔线离两侧关键样本更远，同时惩罚分错或离边界太近的点。软间隔允许一定违例，用\(C\)权衡间隔与这些代价。


\textbf{Gaussian NB先估计每类的数据分布。} 学出各类均值、方差与先验后，比较两类对新样本提供的证据。共享方差使二次项抵消，最后也能整理成线性分数；但这条线来自分布比较。


哪条线最好取决于数据是否适合相应假设、样本数量及正则强度，不能只凭“都线性”判断一样，也不能给出固定冠军。用训练内验证比较候选方案，再用独立测试集评价。


选读：LR最大化的是标签条件似然\(\prod_i P(y_i\mid x_i;w,b)\)。这里的“后验概率”若指\(P(y\mid x)\)，随机对象是类别；若讨论参数MAP，则是给定数据后的\(p(w\mid D)\)。写清对象可以避免把两种后验混为一谈。


\textbf{出处：} 2021A期中 Q8，原卷第3页（10分）


\textbf{答案依据：} 2021A配套解答 / 2021A配套解答副本 Q8随卷解答；不把LR标签似然称为参数后验。
